% What graphifying the harness changes: the released representation (left),
% the typed composition graph that replaces it (middle), and the unfolded
% execution graph one invocation of it writes (right). Input from a figure
% environment; needs \usetikzlibrary{positioning,arrows.meta,fit,calc}.
\begin{tikzpicture}[
  font=\footnotesize,
  panel/.style={draw, rounded corners=2pt, inner sep=7pt},
  code/.style={font=\ttfamily\scriptsize, align=left, anchor=north west},
  gnode/.style={draw, rounded corners=1pt, minimum height=4.4mm, inner xsep=4pt,
                font=\scriptsize, fill=white},
  slot/.style={draw, rounded corners=4pt, inner sep=2.5pt, font=\tiny, fill=gray!8},
  ev/.style={draw, circle, minimum size=5.4mm, inner sep=0pt, font=\tiny, fill=white},
  cone/.style={ev, fill=gray!20},
  dec/.style={-{Latex[length=1.5mm]}, thin},
  obs/.style={-{Latex[length=1.5mm]}, thin, densely dashed},
  wide/.style={-{Latex[length=2.4mm]}, thick},
  note/.style={font=\scriptsize, align=left, text width=4.1cm, anchor=north west,
               inner sep=0pt}]

% ---------------- left: the released representation ----------------------
\node[code] (yaml) at (0,0) {processors:\\
  \ \ - compaction\\
  \ \ - cost\_guard\\
  \ \ - memory.sliding\\
  tools: [Bash, Read, \ldots]\\
  prompt: |\\
  \ \ You are an agent\ldots};
\node[panel, fit=(yaml), label={[font=\scriptsize\bfseries]above:config.yaml}] (L) {};

% ---------------- middle: the typed composition graph -------------------
\node[gnode, fill=gray!12] (h1) at (8.6,-0.1) {hook: before\_model};
\node[gnode, below=5.5mm of h1] (p1) {compaction};
\node[gnode, below=5.5mm of p1] (p2) {cost\_guard};
\draw[dec] (p1) -- node[right, font=\tiny, inner sep=1pt] {attaches} (h1);
\draw[dec] (p2.east) to[bend right] node[right, font=\tiny, inner sep=1pt] (attachlabel) {attaches} (h1.east);
\draw[dec] (p1) -- node[right, font=\tiny, inner sep=1pt] {orders} (p2);
\node[panel, fit=(h1)(p1)(p2)(attachlabel),
      label={[font=\scriptsize\bfseries]above:composition graph $G$}] (M) {};

% ---------------- right: the unfolded execution graph -------------------
\node[cone] (e1) at (15.6,-0.35) {$q_1$};
\node[cone, below right=3.4mm and 5mm of e1] (e2) {$t_1$};
\node[cone, below left=3.4mm and 5mm of e2] (e3) {$q_2$};
\node[ev, below=3.8mm of e3] (e4) {fail};
\draw[dec] (e1) -- (e2);
\draw[obs] (e2) -- node[right, font=\tiny, inner sep=1pt] {data} (e3);
\draw[dec] (e3) -- (e4);
\node[panel, fit=(e1)(e2)(e3)(e4),
      label={[font=\scriptsize\bfseries]above:execution graph $U$}] (Rp) {};

% ---------------- the two transformations -------------------------------
\draw[wide] (L.east) -- node[above, font=\scriptsize] {re-expressed} (M.west);
\draw[wide] (M.east) -- node[above, font=\scriptsize] {one run} (Rp.west);

% ---------------- notes, on one baseline --------------------------------
\coordinate (base) at ($(Rp.south)+(0,-0.32)$);
\node[note] (n1) at (L.west |- base) {\textbf{Released.} Processor order is
  list order plus registration side effects, and an edit is prose: the release has no
  downstream identity check between the candidate description and the
  applied change.};
\node[note] (n2) at (M.west |- base) {\textbf{Graphified.} The schema
  defines six node types; this exporter instantiates hooks, processors and
  slots. $G$ orders graphed processors at byte parity; \texttt{UNGRAPHED}
  processors remain explicit. Hashes identify the graph surface.};
\node[note] (n3) at (Rp.west |- base) {\textbf{Per invocation.} Tool and
  nested-agent events appear in $U$; graphed processors retain their shared
  identifiers. Shaded: the failure's ancestor cone. Unsupported target
  claims are refused by the sixth gate.};
\end{tikzpicture}
